\section{Mecanismo de Attention: la atención es todo lo que necesitas}

\subsection{Más allá de la recurrencia: en busca de un nuevo paradigma}

\begin{figure}[H]
\centering
\includegraphics[width=0.85\textwidth]{slides-images/slide-60.jpg}
\caption{El objetivo del modelado de secuencias: de la secuencia de entrada a las características y de ahí a la salida\protect\footnotemark}
\end{figure}
\footnotetext{Diapositiva 60.}

Volvamos al objetivo fundamental del modelado de secuencias: dada una secuencia de entrada, extraer características y generar una predicción. Las RNN capturan las dependencias de la secuencia procesándola paso a paso, pero ¿podríamos hacerlo mejor si pudiéramos \textbf{ver toda la secuencia a la vez} e identificar automáticamente las partes importantes?

Una idea ingenua sería concatenar todas las entradas en un solo vector grande y procesarlo con una red feedforward. Pero esto presenta dos problemas:
\begin{itemize}
    \item El número de parámetros crece de manera explosiva con la longitud de la secuencia, lo cual no es escalable
    \item La operación de concatenación destruye la información de posición y de orden
\end{itemize}

\begin{figure}[H]
\centering
\includegraphics[width=0.75\textwidth]{slides-images/slide-62.jpg}
\caption{El problema del método ingenuo: concatenar las entradas pierde la información de posición\protect\footnotemark}
\end{figure}
\footnotetext{Diapositiva 62.}

\subsection{Qué es Attention}

La idea central del mecanismo de \textbf{Attention} (atención) proviene de la cognición humana: cuando observamos una escena, no prestamos atención de manera uniforme a todo su contenido, sino que nos enfocamos automáticamente en las partes más relevantes e importantes.

\begin{importantbox}{Definición central de Attention}
Attention es un mecanismo que permite a la red \textbf{identificar automáticamente y enfocarse en las partes más importantes de la entrada}. Calcula \textbf{puntuaciones de relevancia} (Attention Score) entre las distintas partes de la entrada, para decidir a cuáles de ellas debe «prestar atención» al generar la salida.
\end{importantbox}

\subsection{Analogía con un motor de búsqueda: Query-Key-Value}

Quien dio la clase usa una analogía muy intuitiva con un motor de búsqueda para explicar el funcionamiento de Attention:

\begin{figure}[H]
\centering
\includegraphics[width=0.75\textwidth]{slides-images/slide-67.jpg}
\caption{Analogía con un motor de búsqueda: el marco Query-Key-Value\protect\footnotemark}
\end{figure}
\footnotetext{Diapositiva 67.}

Supongamos que buscas «deep learning course» (Query) en un motor de búsqueda de videos:
\begin{enumerate}
    \item Cada video en la base de datos tiene una etiqueta o descripción (Key)
    \item El motor de búsqueda calcula la \textbf{similitud} entre la Query y cada Key
    \item Una vez encontrada la Key que mejor coincide, se extrae el contenido de video correspondiente (Value)
\end{enumerate}

\begin{knowledgebox}{La terna Query-Key-Value}
Los tres conceptos centrales del mecanismo de Attention:
\begin{itemize}
    \item \textbf{Query (Q)}: lo que estás buscando, la «pregunta» que necesita responderse en ese momento
    \item \textbf{Key (K)}: la «etiqueta» de cada entrada en la base de datos, el índice que se usa para hacer coincidir
    \item \textbf{Value (V)}: el contenido real asociado a cada Key, la información que se extrae una vez lograda la coincidencia
\end{itemize}
El proceso de Attention consiste en usar la Query para buscar entre las Keys la entrada que mejor coincide, y luego extraer los Values correspondientes.
\end{knowledgebox}

\subsection{Resumen de la sección}

El núcleo del mecanismo de Attention es un proceso de «búsqueda y recuperación»: se calcula la similitud entre la Query y las Keys para determinar a qué partes de la entrada se debe prestar atención, y luego se extrae la información relevante a partir de los Values. Este mecanismo no necesita procesar la secuencia paso a paso y puede modelar directamente la dependencia entre dos posiciones cualesquiera.

\newpage
